\begin{table}[!h]
\centering
\caption{\label{tab:state_local_totals}Women's Representation: State vs Local Government}
\centering
\fontsize{8}{10}\selectfont
\begin{threeparttable}
\begin{tabular}[t]{lrrr}
\toprule
Level of Government & Total Representatives & Women & \% Women\\
\midrule
\addlinespace[0.3em]
\multicolumn{4}{l}{\textit{\textbf{State Legislatures (1950--2024)}}}\\
\hspace{1em}All State Legislators & $\approx$ 40,000 & $\approx$ 2,000 & 5.0\\
\addlinespace[0.3em]
\multicolumn{4}{l}{\textit{\textbf{Local Bodies (Current)}}}\\
\hspace{1em}Gram Panchayats & $\approx$ 250,000 & --- & ---\\
\hspace{1em}Total Elected Representatives & $>$ 2,000,000 & $\approx$ 1,000,000 & $\approx$ 50.0\\
\bottomrule
\end{tabular}
\begin{tablenotes}
\item \textit{Note: } 
\item State legislature data compiled from historical records of all state elections since 1950. Local body data from the Ministry of Panchayati Raj and Local Government Directory (2024). The contrast illustrates how gender quotas have dramatically increased women's descriptive representation at the local level compared to higher levels where quotas are not mandated.
\end{tablenotes}
\end{threeparttable}
\end{table}
